%%%%%%%%%%%%%%%%%%%%%%%%%%%%%%%%%
%
%       EXERCÍCIO
%
%%%%%%%%%%%%%%%%%%%%%%%%%%%%%%%%%

\definevalorquestao

\begin{exercicioBanco}[\valorquestao]
No mercado de segunda mão, considera-se que determinado \textit{smartphone} perde, a cada ano de uso, 
$\dfrac{1}{4}$ de seu valor em relação ao ano anterior. 

Um consumidor comprou um aparelho por R\$ \(\n{2560,00}\) e, após 4 anos de uso, decidiu substituí-lo por um modelo novo. 
Durante esse período, o preço do modelo novo que ele pretende adquirir permaneceu constante.

Para realizar a troca, o consumidor pode optar por uma das seguintes alternativas:

\begin{enumerate}[(1)]
    \item vender o aparelho usado no mercado de segunda mão, pelo seu valor após 4 anos de depreciação, e utilizar o dinheiro obtido para comprar o novo modelo;
    \item entregar o aparelho usado à loja onde comprará o novo modelo, recebendo um desconto correspondente a $30\%$ do preço original de compra do aparelho.
\end{enumerate}

Considerando apenas o aspecto financeiro, determine qual das duas alternativas é mais vantajosa para o consumidor. (Dica: \(\n{2560} = 2^{9}\cdot 5\)). 

% Define as alternativas
\newcommand{\alternativas}{%
    \begin{center}
        \begin{tabularx}{\textwidth}{X}
            (a) Vender no mercado de segunda mão, pois obterá R\$ 42,00 a mais do que na oferta da loja. \\[3pt]
            (b) Entregar o aparelho na loja, pois obterá R\$ 42,00 a mais do que no mercado de segunda mão. \\[3pt]
            (c) Vender no mercado de segunda mão, pois obterá R\$ 170,00 a mais do que na oferta da loja. \\[3pt]
            (d) Entregar o aparelho na loja, pois obterá R\$ 100,00 a mais do que no mercado de segunda mão. \\[3pt]
            (e) Ambas as opções resultam exatamente no mesmo valor financeiro.
        \end{tabularx}
    \end{center}
}

% Define a resposta correta
\newcommand{\resposta}{A}

% Lógica condicional para exibição
\ifdefstring{\atividade}{lista}{%
    \alternativas
    \vspace{0.5em}
    
    \noindent\textbf{Resposta:} letra \textbf{\resposta}.
}{%
    \ifdefstring{\modo}{objetivo}{%
        \alternativas
    }{%
        % modo = discursiva → não mostra alternativas
    }
}
\end{exercicioBanco}

